\documentclass[11pt,a4paper]{article}
\usepackage[T2A]{fontenc}
\usepackage[utf8]{inputenc}
\usepackage[russian,english]{babel}
\usepackage{amsmath,amssymb,amsthm,mathtools}
\usepackage{setspace}
\usepackage{enumitem}
\usepackage[hidelinks]{hyperref}
\setstretch{1.3}
\frenchspacing
\hyphenpenalty=10000

\newcommand{\head}[1]{\par\noindent\textbf{\underline{#1}}\quad}
\newcommand{\sheettitle}[4]{% title, subtitle, homework-flag, date
  \begin{center}
  {\huge\bfseries #1\par}\vspace{\dimexpr1.2em-1pt\relax minus 3.6pt}
  {\Large #2\par}\vspace{\dimexpr1em+1pt\relax minus 1.2pt}
  \if\relax\detokenize{#3}\relax\else{\Large #3\par}\vspace{\dimexpr1.15em-1pt\relax minus 3.6pt}\fi
  {\foreignlanguage{russian}{#4}\par}
  \end{center}\vspace{0.5pt minus 1.2pt}}
\newcommand{\src}[1]{\unskip\nobreak\hfil\penalty50\hskip1em\hbox{}\nobreak\hfill(#1){\parfillskip=0pt\par}}
\DeclareMathOperator{\tr}{tr}
\DeclareMathOperator{\adj}{adj}
\newcommand{\T}{\mathsf T}
\DeclareMathOperator{\rank}{rank}
\begin{document}
\sheettitle{Linear algebra -- 12}{Mixed problems: commutators, traces and the numerical range}{Homework}{12 сентября 2026 года}

\head{Theoretical minimum.} The trace form $\tr(XY^\T)$ turns the space of matrices into a
Euclidean space, and a hypothesis quantified over all $X$ is a statement about orthogonality in
it. A matrix is a commutator exactly when its trace vanishes, so taking traces is the standard
way to destroy commutators; conversely an identity involving $AB-BA$ usually forces nilpotency
rather than vanishing. Schur's theorem triangularises unitarily and is the right opening move
whenever a spectral hypothesis meets an inequality. The numerical range is convex and contains
the spectrum. For real matrices, non-real eigenvalues come in conjugate pairs, which is why
determinants of real matrices satisfying algebraic identities are so often forced to be
non-negative. And when nothing else works: apply the identity to an eigenvector, or restrict to a
$2\times2$ block, or read the whole hypothesis as a linear map on $M_n$.

\par\vspace{-1pt}\noindent\rule{\linewidth}{0.4pt}\par\vspace{7pt}

\begin{enumerate}
\item (AKHIMO 2018) Let $n\geqslant2$ be an integer and denote by $A_1,\dots,A_n\in M_2(\mathbb C)$
all solutions of the equation $X^n=\begin{pmatrix}2&3\\4&6\end{pmatrix}$. Prove that
$\sum_{k=1}^n\tr A_k=0$.

\item (AKHIMO 2022) Let $A,B,C\in M_n(\mathbb R)$ be skew-symmetric matrices with
$\det(A^2+B^2+C^2)=0$. Prove that $\det(AP+BQ+CR)=0$ for every triple
$P,Q,R\in M_n(\mathbb R)$.

\item (SEEMOUS 2025) Let $A$ be an $n\times n$ matrix with strictly positive entries and let
$u,v\in\mathbb R^n$ have strictly positive entries with $Au=v$ and $Av=u$. Prove that $u=v$.

\item (SEEMOUS 2024) Let $A,B\in M_n(\mathbb R)$ be symmetric with non-negative eigenvalues. Prove
that $A^3+B^3=(A+B)^3$ if and only if $AB=O_n$.

\item (SEEMOUS 2023) Prove that if $A$ and $B$ are complex $n\times n$ matrices satisfying
\[
A=AB-BA+A^2B-2ABA+BA^2+A^2BA-ABA^2,
\]
then $\det A=0$.

\item (SEEMOUS 2021) Let $n\geqslant2$ and let $A\in M_n(\mathbb R)$ satisfy $A^2=-E$. If
$B\in M_n(\mathbb R)$ and $AB=BA$, prove that $\det B\geqslant0$.

\item (SEEMOUS 2017) Let $A=\begin{pmatrix}a&b\\c&d\end{pmatrix}\in M_2(\mathbb R)$ satisfy
$a^2+b^2+c^2+d^2<\frac15$. Show that $E+A$ is invertible.

\item (SEEMOUS 2016) Let $A_1,\dots,A_k\in M_n(\mathbb R)$ be idempotent. Prove that
\[
\sum_{i=1}^kN(A_i)\geqslant\rank\Bigl(E-\prod_{i=1}^kA_i\Bigr),\qquad\text{where }
N(A_i)=n-\rank A_i.
\]

\item (AKHIMO 2022) Let $A=\operatorname{diag}(\lambda_1,\dots,\lambda_n)$ with real $\lambda_i$.
Prove that there exist a unitary matrix $U$ and a real number $\lambda$ such that
$|AU-UA|\geqslant|A-\lambda E|$, where $|B|=\sqrt{B^*B}$ and $X\geqslant Y$ means that $X-Y$ is
positive semidefinite.

\item (SEEMOUS 2015) For an integer $n>2$ let $A,B,C,D\in M_n(\mathbb R)$ satisfy $AC-BD=E$ and
$AD+BC=O$. Prove that
\begin{enumerate}[label=\alph*),align=left]
\item $CA-DB=E$ and $DA+CB=O$;
\item $\det(AC)\geqslant0$ and $(-1)^n\det(BD)\geqslant0$.
\end{enumerate}

\item (IMC) Let $A,B\in M_n(\mathbb C)$ satisfy $A^2B+BA^2=2ABA$. Prove that $(AB-BA)^k=O$ for
some positive integer $k$.

\item (VJIMC 2010) Let $A,B$ be complex $2\times2$ matrices with $AB-BA=B^2$. Prove that $AB=BA$.

\item (IMC) Let $A,B$ be real $n\times n$ matrices with $\rank(AB-BA+E)=1$. Prove that
$\tr(ABAB)-\tr(A^2B^2)=\frac12n(n-1)$.

\item (IMC) Let $A,B,C$ be real square matrices of the same size with $A$ invertible. Prove that
if $(A-B)C=BA^{-1}$, then $C(A-B)=A^{-1}B$.

\item Prove the Toeplitz--Hausdorff theorem: the numerical range $W(A)$ is convex. Deduce its real
version: for real symmetric $A,B$ and $n\geqslant3$ the set
$\{(\langle Ax,x\rangle,\langle Bx,x\rangle):\|x\|=1\}$ is convex.

\item (after IMC 2025 shortlist) Let $A,B$ be real $n\times n$ matrices with
$\tr\bigl((AB)^2\bigr)=\tr(A^2B^2)$.
\begin{enumerate}[label=\alph*),align=left]
\item Prove that if all eigenvalues of $AB-BA$ are real, then $AB-BA$ is nilpotent.
\item Show that for $n=2$ the assumption on the eigenvalues can be dropped, but for every
$n\geqslant3$ it cannot.
\end{enumerate}
\end{enumerate}
\end{document}
